\begin{notechapter}
\chapter{Vector Spaces and Duality}
\label{ch:duality}
\chaptersubtitle{%
  Covectors, dual bases and the dual operator}

\begin{chapterplan}
\begin{itemize}[nosep]
  \item Vector spaces, bases and dimension; linear maps and their matrices;
        change of basis.
  \item The dual space $V^*$: linear forms and their level sets; the dual basis;
        how the components of vectors and of covectors change with the basis.
  \item The dual operator $\dual A:W^*\to V^*$, $\phi\mapsto\phi\circ A$. It
        needs no inner product, and its matrix in the dual bases is the
        transpose of the matrix of $A$, over any field and with no conjugation.
  \item The bidual and the canonical isomorphism $V\cong V^{**}$; why an
        isomorphism $V\cong V^*$ always involves a choice.
  \item Annihilators; $\ker\dual A = (\im A)^\circ$; $A$ and $\dual A$ have the
        same rank.
  \item The basis test: which formulas survive a change of basis.
\end{itemize}
\smallskip\noindent\textbf{Machine learning threads.}
\begin{itemize}[nosep]
  \item Row vectors and column vectors.
  \item Coordinatewise optimizers such as Adam depend on the basis.
  \item The vector--Jacobian product is a dual operator
        (\cref{ch:chain-rule-autodiff}).
\end{itemize}
\end{chapterplan}
\end{notechapter}
